%%%%%%%%%%%%%%%%%%%%%%%%
%
% $Autor: Wings $
% $Project:  $
% $Filename: OledSSD1306
% $Short description: In diesem Abschnitt wird das OLED-Display genauer beschrieben. Wie dieses aufgebaut ist und was benötigt wird, um dieses in ein Projekt einzubinden.
%
%%%%%%%%%%%%

\chapter{OLED-Display}\label{sec:OLEDDisplay}
\section{Allgemein}
\ac{oled}s (Organic Light Emitting Diodes) stellen eine besondere Form der LED-Technologie dar, bei der jedes Pixel aus eigenständig leuchtenden organischen Leuchtdioden besteht. Im Gegensatz zu LCDs mit LED-Hintergrundbeleuchtung benötigen OLEDs keine Hintergrundbeleuchtung oder LC-Zellen, da die einzelnen Pixel selbst Licht erzeugen. Dies ermöglicht nicht nur eine besonders flache Bauweise, sondern auch eine exakte Steuerung einzelner Bildpunkte – inklusive vollständigem Abschalten für tiefstes Schwarz und hohen Kontrast.

Ein wesentlicher Vorteil gegenüber LCD-Technologien ist die deutlich schnellere Reaktionszeit: Helligkeitsänderungen erfolgen bei OLEDs in unter einer Mikrosekunde. Auch der Farbraum und Kontrastumfang sind bei OLED-Bildschirmen in der Regel erweitert. Allerdings ist die Lebensdauer einzelner OLEDs im Vergleich zu anderen Bildschirmtechnologien tendenziell geringer, obwohl bereits von Betriebszeiten im Bereich von 15 Jahren gesprochen wird. 
In der Praxis kommen verschiedene OLED-Technologien zum Einsatz, insbesondere RGB-SBS-AMOLED (mit drei nebeneinander angeordneten Farb-OLEDs pro Pixel) und WOLED (weiße OLEDs mit nachgeschalteten Farbfiltern). OLED-Elemente werden meist im Dünnschichtverfahren hergestellt, was die Produktion vereinfacht.\cite{Stotz.2019}

Für Projekte mit Mikrocontrollern wie den Arduino oder Rasperry Pi werden oft monochrome OLED-Displays verwendet. Das bedeutet, dass anstelle von RGB-LEDs nur normale, einfarbige LEDs verwendet werden. Oft sind diese monochromen OLED-Displays blau oder weiß, um einen möglichst hohen Kontrast zu bieten. \cite{Schreiter.2019}



\section{OLED-Display mit SSD1306 Controller}
Das \ac{oled}-Display mit dem Controller SSH1106 dient der Ausgabe von Messwerten des Arduinos. Es besitzt Abmessungen von 36 × 34 × 3 mm und überzeugt durch seinen hohen Kontrast sowie gute Lesbarkeit. Das Display verfügt über eine Punktmatrix von 128 × 64 Pixeln, wobei jedes Pixel eine Größe von 0,21 mm × 0,21 mm bei einem Abstand von 0,23 mm × 0,23 mm aufweist. Die Darstellung erfolgt in weißer Farbe im passiven Matrixmodus. Angesteuert wird das Display über die I\textsuperscript{2}C-Schnittstelle (Pins VCC, GND, SCL, SDA) und benötigt eine Betriebsspannung von 3,3 V bis 5 V bei einem Stromverbrauch unter 11 mA. Der zulässige Betriebstemperaturbereich liegt zwischen -20 °C und +70 °C. Für die Ansteuerung kommt die U8glib-Bibliothek zum Einsatz, welche eine Vielzahl von Grafikfunktionen für \ac{oled}-Displays mit SSH1106-Controller bereitstellt.\cite{Schreiter.2019}

\begin{longtable}[h]{|p{5cm}|p{8cm}|}
	\caption{Spezifikationen OLED-Display}
	\label{tab:specs_oled}\\
	\hline
	
	% Hauptmerkmale
	\multicolumn{2}{|l|}{\rule{0pt}{15pt}\textbf{Hauptmerkmale}} \\ \hline
	Modell & OLED-Display \\ \hline
	Besonderheit & Kleines Display mit hoher Auflösung, leicht anzuschließen, universell \\ \hline
	Abmessungen & 28 x 33 x 4 \\ \hline
	Gewicht & 4 g \\ \hline
	Lieferumfang & SBC-OLED01 \\ \hline
	
	% Weitere Spezifikationen
	\multicolumn{2}{|l|}{\rule{0pt}{15pt}\textbf{Weitere Spezifikationen}} \\ \hline
	Spannungsversorgung & 3 V bis 5 V \\ \hline
	Auflösung & 128 x 64 Pixel \\ \hline
	Pixelgröße & 0{,}152 x 0{,}152 mm \\ \hline
	Hintergrundfarbe & Schwarz \\ \hline
	Anzeigefarbe & Weiß \\ \hline
	Anschluss & 4-polig, 2{,}54 mm Raster \\ \hline
	Interface & I\textsuperscript{2}C \\ \hline
	Controller-Chip & SSD1306 \\ \hline
	Stromverbrauch & {<}11 mA \\ \hline
	
\end{longtable}


\subsection{Aufbau und Funktionsweise}
In der Abbildung \ref{fig:OLEDsketch} wird der Aufbau des OLED-Displays gezeigt. Über die vier Pins kann später die Stromversorgung und Logikspannung des OLED-Displays mit dem Mikrocontroller verbunden werden.

\begin{center}
\begin{circuitikz}
    \centering
    \tikzstyle{every node}=[font=\Large]
    \draw [ fill={rgb,255:red,24; green,9; blue,78} , line width=1pt ] (2.5,18.5) rectangle (11.25,7.25);
    
    \node [font=\scriptsize, color={rgb,255:red,255; green,255; blue,255}, rotate around={90:(0,0)}] at (3.5,13.25) {SCK};
    \draw [ fill={rgb,255:red,0; green,0; blue,0} , line width=1pt ] (4,18) rectangle (9.75,7.75);
    \draw [ fill={rgb,255:red,40; green,10; blue,123} , line width=1pt ] (4.25,17.75) rectangle (9,8);
    \draw [ fill={rgb,255:red,188; green,184; blue,184} ] (10.5,17.75) circle (0.5cm);
    \draw [ dashed] (10.5,17.75) circle (0.25cm);
    \draw [ fill={rgb,255:red,188; green,184; blue,184} ] (3.25,8) circle (0.5cm);
    \draw [ fill={rgb,255:red,188; green,184; blue,184} ] (3.25,17.75) circle (0.5cm);
    \draw [ fill={rgb,255:red,188; green,184; blue,184} ] (10.5,8) circle (0.5cm);
    \node [font=\LARGE] at (19.25,10.5) {};
    \draw [ dashed] (10.5,8) circle (0.25cm);
    \draw [ dashed] (3.25,17.75) circle (0.25cm);
    \draw [ dashed] (3.25,8) circle (0.25cm);
    \draw [ fill={rgb,255:red,186; green,186; blue,186} , line width=0.2pt , rounded corners = 9.6] (10,14.75) rectangle (11.25,11);
    \draw [ color={rgb,255:red,186; green,186; blue,186} , fill={rgb,255:red,186; green,186; blue,186}, line width=0.2pt ] (10.75,14.75) rectangle (11.25,11);
    \draw [ color={rgb,255:red,255; green,0; blue,0} , fill={rgb,255:red,255; green,184; blue,184}, dashed] (2.75,14.25) rectangle  (3.25,13.75);
    \draw [ color={rgb,255:red,255; green,0; blue,0} , fill={rgb,255:red,255; green,184; blue,184}, dashed] (2.75,13.5) rectangle  (3.25,13);
    \draw [ color={rgb,255:red,255; green,0; blue,0} , fill={rgb,255:red,255; green,184; blue,184}, dashed] (2.75,12.75) rectangle  (3.25,12.25);
    \draw [ color={rgb,255:red,255; green,0; blue,0} , fill={rgb,255:red,255; green,184; blue,184}, dashed] (2.75,12) rectangle  (3.25,11.5);
    \draw [ color={rgb,255:red,255; green,0; blue,0} , fill={rgb,255:red,255; green,184; blue,184}] (3,14) circle (0.25cm);
    \draw [ color={rgb,255:red,255; green,0; blue,0} , fill={rgb,255:red,255; green,184; blue,184}] (3,13.25) circle (0.25cm);
    \draw [ color={rgb,255:red,255; green,0; blue,0} , fill={rgb,255:red,255; green,184; blue,184}] (3,12.5) circle (0.25cm);
    \draw [ color={rgb,255:red,255; green,0; blue,0} , fill={rgb,255:red,255; green,184; blue,184}] (3,11.75) circle (0.25cm);
    \draw [ color={rgb,255:red,208; green,184; blue,27} , fill={rgb,255:red,208; green,184; blue,27}] (10,14.25) rectangle (10.5,11.5);
    
    %Beschriftung auf Platine
    \node [font=\scriptsize, color={rgb,255:red,255; green,255; blue,255}, rotate around={90:(0,0)}] at (3.5,14) {SDA};
    \node [font=\scriptsize, color={rgb,255:red,255; green,255; blue,255}, rotate around={90:(0,0)}] at (3.5,13.25) {SCL};	
    \node [font=\scriptsize, color={rgb,255:red,255; green,255; blue,255}, rotate around={90:(0,0)}] at (3.5,12.5) {VDD};
    \node [font=\scriptsize, color={rgb,255:red,255; green,255; blue,255}, rotate around={90:(0,0)}] at (3.5,11.75) {GND};
    
    
    \node [font=\large, color={rgb,255:red,255; green,255; blue,255}, rotate around={90:(0,0)}] at (3,14.5) {4};
    \node [font=\large, color={rgb,255:red,255; green,255; blue,255}, rotate around={90:(0,0)}] at (3,11.25) {1};
    \draw [ fill={rgb,255:red,130; green,130; blue,130} ] (10,14) rectangle (10.25,11.75);
    
    %Beschriftung
    \node [font=\Large] at (-1,14) {I²C Serial Data Line - SDA};
    \node [font=\Large] at (-1,13.25) {I²C Serial Clock Line - SCL};
    \node [font=\Large] at (-1,12.5) {Power Supply - VDD};
    \node [font=\Large] at (-1,11.75) {Ground - GND};Ground - GND
    
    \draw[-, thick] (12,10) -- (8,10);
    \node[right] at (12,10) {OLED-Display 128 x 64 Pixel};
\end{circuitikz}
	\captionof{figure}[OLED-Display]{OLED-Display nach \href{../../../Documents/Datenblaetter/DatenblattOLEDDisplay}{[Datenblatt.PDF]}}
	\label{fig:OLEDsketch}
\end{center}

Bei der Ansteuerung eines OLED-Displays über das I\textsuperscript{2}C-Protokoll übernehmen die beiden Leitungen \ac{sda} und \ac{scl} den seriellen Datenaustausch zwischen Mikrocontroller (Master) und OLED-Display (Slave).

\begin{itemize}
	\item SDA überträgt die eigentlichen Datenbits (z.B. Befehle oder Bildinhalte).
	\item SCL gibt den Takt vor, der bestimmt, wann ein Bit gelesen oder geschrieben wird.
\end{itemize}

Das I\textsuperscript{2}C-Protokoll ermöglicht eine einfache, bidirektionale Kommunikation über nur zwei Leitungen, wobei mehrere Geräte über individuelle Adressen angesprochen werden können.

\subsubsection{Protokoll I\textsuperscript{2}C}
In den 1980-er Jahre wurde die I\textsuperscript{2}C-Schnittstelle entwickelt zum Datenaustausch zwischen verschiedenen Komponenten 
eines Gerätes. Im Gegensatz zur seriellen Verbindung, welche stets nur zwei Kommunikationspartner verbindet, handelt es sich hierbei um einen Datenbus, welcher optional deutlich mehr Komponenten (ursprünglich bis zu 112, in neueren Versionen nochmals deutlich mehr) miteinander verbinden kann. 
Elektrisch gesehen besteht der Bus aus lediglich zwei Leitungen, an die alle Busteilnehmer angeschlossen sind, sowie einem gemeinsamen Massebezug. Eine der Leitungen dient als Taktsignal (Clock - CLK), die andere wird zum Datenaustausch (Serial Data – SDA) verwendet. Dabei fungiert stets ein Teilnehmer als Master, er generiert das Taktsignal und koordiniert die Kommunikation – etwa so wie ein Moderator in einer Gesprächsrunde. Alle anderen Komponenten agieren als Slaves, sie ordnen sich also dem Master unter. Jeder Slave benötigt eine eindeutige Adresse (in Form einer Zahl), über die er vom Master angesprochen werden kann. Der Master kann den Slaves jederzeit Daten senden.

\begin{center}
\begin{circuitikz}
    \tikzstyle{every node}=[font=\LARGE]
    \node [font=\LARGE] at (19.25,17.75) {};
    \draw [ line width=0.6pt ] (2.5,23.25) rectangle  node {\LARGE Slave} (6.25,20.75);
    \draw [ line width=0.6pt ] (7.5,23.25) rectangle  node {\LARGE Slave} (11.25,20.75);
    \draw [ line width=0.6pt ] (12.5,23.25) rectangle  node {\LARGE Slave} (16.25,20.75);
    \draw [ line width=0.6pt ] (17.5,23.25) rectangle  node {\LARGE Slave} (21.25,20.75);
    \draw [ line width=0.6pt ] (2.5,17) rectangle  node {\LARGE Slave} (6.25,14.5);
    \draw [ line width=1.2pt ] (7.5,17) rectangle  node {\LARGE Master} (11.25,14.5);
    \draw [ line width=0.6pt ] (12.5,17) rectangle  node {\LARGE Slave} (16.25,14.5);
    \draw [ line width=0.6pt ] (17.5,17) rectangle  node {\LARGE Slave} (21.25,14.5);
    \draw [ line width=0.6pt](3.75,20.75) to[short] (3.75,17);
    \draw [ line width=0.6pt](3.75,19.5) to[short] (18.75,19.5);
    \draw [ line width=0.6pt](5,20.75) to[short] (5,17);
    \draw [ line width=0.6pt](5,18.25) to[short] (20,18.25);
    \draw [ line width=0.6pt](20,20.75) to[short] (20,17);
    \draw [ line width=0.6pt](18.75,19.5) to[short] (18.75,20.75);
    \draw [ line width=0.6pt](13.75,20.75) to[short] (13.75,17);
    \draw [ line width=0.6pt](15,20.75) to[short] (15,17);
    \draw (10,20.75) to[short] (10,17);
    \draw [ line width=0.6pt](8.75,20.75) to[short] (8.75,17);
    \node at (3.75,19.5) [circ] {};
    \node at (5,18.25) [circ] {};
    \node at (8.75,19.5) [circ] {};
    \node at (10,18.25) [circ] {};
    \node at (13.75,19.5) [circ] {};
    \node at (15,18.25) [circ] {};
    \node at (18.75,19.5) [circ] {};
    \node at (20,17) [circ] {};
    \draw [ line width=0.6pt](18.75,19.5) to[short] (18.75,17);
    \node [font=\LARGE] at (12,19.75) {SDA};
    \node [font=\LARGE] at (12,18.5) {SCL};
    \node at (20,18.25) [circ] {};
\end{circuitikz}

	\captionof{figure}{I\textsuperscript{2}C-Bus mit 8 Teilnehmern nach \cite{Schreiter.2019}}
	\label{fig:I2CSketch}
\end{center}

Die einzelnen Slaves dürfen jedoch nur senden, wenn sie vom Master abgefragt werden, also quasi eine Sendeerlaubnis erhalten. Dadurch wird vermieden, dass mehrere Komponenten gleichzeitig auf die einzige vorhandene Datenleitung senden, denn dann wäre das Signal gestört. Üblicherweise agiert unser Arduino natürlich als Master und kommuniziert mit Komponenten wie Sensoren oder OLED-Displays als Slaves. Allerdings ist es auch möglich, über I\textsuperscript{2}C Verbindungen zwischen mehreren Arduinos herzustellen.


\subsection{Anwendung}
\ac{oled}-Displays werden in vielen Bereichen eingesetzt, in denen energieeffiziente und kontrastreiche Anzeigen benötigt werden. Typische Anwendungen finden sich in Fernsehern, tragbaren Geräten wie Smartwatches und medizinischen Geräten.
Seit etwa 15 Jahren werden LCDs in vielen Anwendungsbereichen von OLED-Displays ersetzt. Durch die Möglichkeit Pixel einzeln anzusteuern und auch komplett abzuschalten, kann der OLED-Display sehr kontraststarke Bilder erzeugen, welches ebenfalls für eine gute Lesbarkeit sorgt.
Durch den geringen Stromverbrauch sind OLED-Displays ebenfalls besonders gut für batteriebetriebene Systeme geeignet.\cite{Schreiter.2019}


\section{Installation}
Zur Nutzung des OLED-Displays muss dieses nach mit dem Arduino nach Tabelle \ref{tab:oled_pinbelegung} angeschlossen werden. Bei der Benutzung des Arduino Shields, können Grove-Stecker verwendet werden. 
\begin{longtable}[h]{|p{5cm}|p{8cm}|}
	\caption{Pinbelegung des OLED-Displays}
	\label{tab:oled_pinbelegung} \\
	\hline
	\textbf{OLED-Display Pin} & \textbf{Arduino Pin} \\ \hline
	SDA & A4 \\ \hline
	SCL & A5 \\ \hline
	VCC & {+}3V3 out \\ \hline
	GND & GND \\ \hline
\end{longtable}




\section{Bibliotheken}
Für die Verwendung des OLED-Displays sind die folgenden aufgelisteten Bibliotheken notwendig \ref{tab:oled_libs}. Um eine Bibliothek in ein Programm zu integrieren muss diese in der Arduino IDE heruntergeladen und installiert sein und anschließend im Programmcode mit dem folgenden Befehl eingebunden werden: \PYTHON{\#include<name.h>} \newline

 
\section{Verwendete Bibliotheken}
\begin{center}
	\captionof{table}{Verwendete Bibliotheken zum Testen des OLED-Displays}
	\label{tab:oled_libs}
	\begin{tabular}{|p{5cm}|p{9cm}|}
			\hline
			\textbf{Bibliothek} & \textbf{Beschreibung} \\
			\hline
			\textbf{Wire.h} &  \textbf{Version:} Arduino-Standardbibliothek \newline
			Funktion: Dient der Kommunikation über den I\textsuperscript{2}C-Bus zwischen Mikrocontroller und OLED-Display. \\
			\hline
			\textbf{Adafruit\_GFX.h} & \textbf{Version:} 1.12.0 \newline
			Grafikbibliothek zur Darstellung von Text, Linien, Rechtecken und weiteren primitiven Elementen auf grafischen Displays. \\
			\hline
			\textbf{Adafruit\_SSD1306.h} & \textbf{Version:} 2.5.13 \newline
			Treiber für SSD1306-basierte OLED-Displays mit Unterstützung für I\textsuperscript{2}C. \\
			\hline
		\end{tabular}
\end{center}

\subsection{Wire.h}
Die Bibliothek Wire.h ist eine Standardbibliothek für Arduino-Plattformen, welche Funktionen und Methoden für die Kommunikation zur Verfügung stellt. Mit Hilfe der Schnittstelle I\textsuperscript{2}C ermöglicht die Bibliothek die Kommunikation zwischen einem Arduino und einem anderen I\textsuperscript{2}C-fähigen Gerät wie z.B. Sensoren oder Displays. I\textsuperscript{2}C ist ein Kommunikationsbus, der im Vergleich zu seriellen Schnittstellen den Vorteil hat, dass er mit mehr als zwei Geräten kommunizieren kann. Ein I\textsuperscript{2}C-Bus benötigt zwei Leitungen: SCL für ein Taktsignal und SDA für Daten.

\subsection{Adafruit\_GFX.h}
Die Adafruit\_GFX-Bibliothek bietet grundlegende Grafikfunktionen für viele verschiedene LCD- und OLED-Displays sowie LED-Matrizen von Adafruit. Sie sorgt dafür, dass Programme leicht zwischen unterschiedlichen Displays angepasst werden können. Neue Funktionen und Verbesserungen gelten automatisch für alle unterstützten Displays.
Die Bibliothek wird immer zusammen mit einer weiteren Bibliothek verwendet, die speziell für das jeweilige Display gedacht ist, für unser spezifisches OLED-Display, welches den Controller-Chip SSD1306 verwendet, muss die Adafruit\_SSD1306.h Bibliothek ebenfalls mit eingebunden werden.

\subsubsection{Adafruit\_SSD1306.h}
Die Adafruit\_SSD1306 Bibliothek ist speziell für OLED-Displays mit dem SSD1306-Controller entwickelt worden. Sie arbeitet zusammen mit der Adafruit\_GFX-Bibliothek und ermöglicht die Ansteuerung und Darstellung von Grafiken und Text auf dem Display.


%\subsection{Kalibrierung}


\section{Codebeispiel}

\subsection{Kurzbeschreibung des Programms}
Dieses Arduino-Testprogramm dient zur Initialisierung und Darstellung einfacher Texte auf einem SSH1106 OLED-Display. Nach der Initialisierung wird die Nachricht \textit{"Hello, world!"} angezeigt. Anschließend wird die Anzeige im Sekundentakt invertiert, um einen Blinke-Effekt zu erzeugen.

%\subsubsection{Verwendete Bibliotheken}





\subsection{Verwendete Funktionen}
\begin{center}
	\captionof{table}{Funktionen im Testprogramm für das OLED-Display.}
	\begin{tabular}{|p{5cm}|p{9cm}|}
		\hline
		\textbf{Funktion} & \textbf{Beschreibung} \\
		\hline
		\PYTHON{Serial.begin()} & Startet die serielle Kommunikation mit 9600 Baud. \\
		\hline
		\PYTHON{display.begin()} & Initialisiert das OLED-Display mit angegebener I\textsuperscript{2}C-Adresse. \\
		\hline
		\PYTHON{display.clearDisplay()} & Löscht alle Inhalte auf dem Display. \\
		\hline
		\PYTHON{display.setTextSize()} & Setzt die Schriftgröße für Textausgabe. \\
		\hline
		\hline
		\PYTHON{display.setCursor()} & Positioniert den Textcursor an einer bestimmten Stelle. \\
		\hline
		\PYTHON{display.println()} & Gibt eine Textzeile auf dem Display aus. \\
		\hline
		\PYTHON{display.display()} & Überträgt die Zeichenpufferdaten auf das Display. \\
		\hline
		\PYTHON{display.invertDisplay()} & Invertiert die Anzeige des Displays. \\
		\hline
		\hline
	\end{tabular}
\end{center}

\subsection{Globale Variablen und Parameter}
\begin{center}
	\captionof{table}{Globale Variablen und Parameter im Testprogramm für das OLED-Display.}
	\begin{tabular}{|p{5cm}|p{9cm}|}
		\hline
		\textbf{Variable} & \textbf{Beschreibung} \\
		\hline
		\PYTHON{SCREEN\_WIDTH} & Definiert die Breite des Displays in Pixeln (128). \\
		\hline
		\PYTHON{SCREEN\_HEIGHT} & Definiert die Höhe des Displays in Pixeln (64). \\
		\hline
		\PYTHON{display} & Objekt zur Steuerung des OLED-Displays mit Adafruit SSD1306-Treiber. \\
		\hline
	\end{tabular}
\end{center}

\subsection{Testcode}
{
	\captionof{code}{Testprogramm für das OLED-Display}
	\label{Code:Test:File:TestOLEDDisplay}
	\ArduinoExternal{}{../../Code/Arduino/OledSSD1306/Test/Test.ino}  
}


\section{Weiterführende Literatur}

\begin{itemize}
	\item 	\textbf{OLED:}
	Weiterführende Informationen zum Aufbau und zur Funktionsweise von OLEDs enthält das Buch 
	\glqq \textit{Elektronik für Ingenieure und Naturwissenschaftler}\grqq \cite{Hering.2021}. 
	
\begin{itemize}
	\item Beschreibt sämtliche Komponenten der Elektrotechnik, welche für Ingenieure relevant sein können.
	\item Sehr ausführliche Beschreibung der Struktur und Aufbau von LEDs und auch speziell OLED-Displays.
	\item Nennt weitere Displaytypen und Alternativen im Vergleich zu OLED.
\end{itemize}
	
	\item 	\textbf{Ansteuerung:} Sehr gute inoffizielle Anleitung zum Basteln mit dem Arduino und viele Anwendungsbeispiele bietet das Buch 	
	\glqq \textit{Arduino-Kompendium}\grqq \cite{Schreiter.2019}. 
	
	\begin{itemize}
		\item Bietet eine Vielzahl von Kurzanleitungen zu vielen Sensoren und Komponenten die teilweise im offiziellen Handbuch von Arduino fehlen.
		\item Enthält viele Praxisprojekte zum Nachbauen
	\end{itemize}
		
	
	\item \textbf{Anzeigemöglichkeiten auf dem OLED-Display}: Um herauszufinden was man alles mit dem OLED-Display machen kann, ist die Dokumentation der Adafruit\_GFX.h Bibliothek zu empfehlen. \href{https://learn.adafruit.com/adafruit-gfx-graphics-library}{learn.adafruit.com/adafruit-gfx-graphics-library}. 
	Auf der Webseite werden sämtliche Funktionen der Bibliothek dokumentiert sowie weitere Bibliotheken verlinkt, die für Funktionen des OLED-Displays kompatibel sind.
	\begin{itemize}
		\item Anzeigen von Bildern, Animationen oder Texten auf dem Display
		\item Laden und nutzen von größeren Dateien in den RAM von externen Speichergeräten wie einer SD-Karte
		\item Weitere Bibliotheken von Adafruit 
	\end{itemize}

\end{itemize}




% Structure
\chapter{1,3 Zoll OLED Display}
\label{ch:Display}

Das 1,3 Zoll OLED Display dient der Visualisierung von Rückmeldungen an den Benutzer. In den Folgenden Unterkapiteln wird auf den allgemeinen Aufbau und die Funktionsweise als auch auf das spezifische Display und dessen Anwendung mit Code eingegangen.

\section{Grundlagen der OLED-Technologie}

\subsection{Was ist eine OLED?}
Eine Organische Leuchtdiode (OLED) ist ein dünnfilmiges Bauelement, das Licht durch elektrische Anregung organischer Halbleitermaterialien emittiert. Im Gegensatz zu herkömmlichen LEDs basieren OLEDs nicht auf anorganischen Kristallen (z.B. Galliumnitrid), sondern auf Kohlenwasserstoffverbindungen \cite{Bartlett:2023} . Dieser Unterschied ermöglicht: 

\begin{itemize}
    \item Flexible Bauformen (z.B. auf Polymerfolien)
    \item Flächige, gleichmäßige Beleuchtung ohne Hintergrundlicht
    \item Hohe Energieeffizienz durch Direktemission
\end{itemize}
\cite{Bartlett:2023} \cite{Stotz:2019} \cite{Wellmann:2017} 

\subsection{Aufbau einer OLED}
Abbildung \ref{fig:OLEDAufbau} zeigt den typischen Schichtaufbau:
\begin{itemize}
    \item Substrat: Trägermaterial aus Glas oder flexibler Kunststofffolie (z.B. Polyimid)
    \item Anode: Transparentes Indium-Zinn-Oxid (ITO) zur Injektion positiver Ladungsträger
    \item Organische Funktionsschichten (Gesamtdicke < 500 nm):
    \begin{itemize}
        \item Lochtransportlayer (HTL): Leitet positive Ladungen (z.B. Kupferphthalocyanin)
        \item Emitterschicht (EML): Erzeugt Licht durch Rekombination (z.B. Alq\textsubscript{3})
        \item Elektronentransportlayer (ETL): Leitet negative Ladungen (z.B. BCP)
    \end{itemize}
    \item Kathode: Metallschicht (z.B. Al/Ca) für Elektroneninjektion
\end{itemize}
\cite{Bartlett:2023} \cite{Stotz:2019} \cite{Wellmann:2017}

\begin{figure}[htbp]
    \centering
    \resizebox{0.8\linewidth}{!}{%
        \input{../../Images/OLED/tikzOLED}
    }
    \caption{Vereinfachte Darstellung des OLED-Schichtaufbaus (adaptiert nach Wellmann, 2017, Kap. 3.4.2 Abb.3-83) \cite{Wellmann:2017}}
    \label{fig:OLEDAufbau}
\end{figure}

\subsection{Funktionsprinzip}
OLEDs nutzen Elektrolumineszenz – die Umwandlung elektrischer Energie in Licht. Dieser Prozess erfolgt in vier Schritten:
\begin{enumerate}
    \item \textbf{Ladungsinjektion:}
    Bei Anlegen einer Gleichspannung (3–5 V) injiziert die Kathode Elektronen in die ETL, während die Anode Löcher (positive Ladungsträger) in die HTL einbringt.
    
    \item \textbf{Ladungstransport:}  
    Die Ladungsträger wandern durch die organischen Schichten – Elektronen zur EML, Löcher zur EML.  
    
    \item \textbf{Rekombination:}  
    In der EML kombinieren Elektronen und Löcher unter Freisetzung von Energie. Nur 20–25\% dieser Energie wird als Licht abgegeben (sichtbarer Bereich: 450–650 nm), der Rest geht als Wärme verloren.  
    
    \item \textbf{Lichtaustritt:}  
    Photonen durchdringen die transparente Anode. Optische Verluste durch Reflexion begrenzen die Effizienz auf ca.20\%  \cite{Wellmann:2017}.  
\end{enumerate}

\subsection{Stärken von OLEDs}
\begin{itemize}
    \item Bildqualität:
    Nahezu unbegrenzter Kontrast (da Pixel komplett abschaltbar)
    \item Bauhöhe:
    Ultraflach (< 1 mm) durch fehlende Hintergrundbeleuchtung
    \item Energieverbrauch:
    Geringer Strombedarf bei dunklen Bildinhalten
\end{itemize}
\cite{Bartlett:2023} \cite{Stotz:2019} \cite{Wellmann:2017}

\subsection{Herausforderungen}
\begin{itemize}
    \item Lebensdauer:
    Organische Materialien degradieren unter Sauerstoff- und Feuchteeinfluss 
    \item Herstellungskosten:
    Hochvakuumprozesse für SM-OLEDs sind teurer als LCD-Fertigung
    \item Helligkeit:
    Begrenzte Maximalhelligkeit verglichen mit Micro-LEDs
\end{itemize}

\cite{Bartlett:2023} \cite{Stotz:2019} \cite{Wellmann:2017}



\section{1,3 Zoll OLED Display}
Das 1,3 Zoll \ac{oled} Display hat eine Auflösung von 128x64 Pixeln. Das Display nutzt die I2C-Schnittstelle für die Kommunikation und hat eine niedrige Stromaufnahme von weniger als 11mA. Durch die kompakte Bauweise und die hohe Effizienz ist dieses Display für verschiedene Anwendungen geeignet. Dieses OLED Display hat den SH1106 Chip verbaut. 

\begin{figure}[h]
    \begin{center}
        \includegraphics[width=12cm]{../../Images/OLED/13OLEDDisplay}
        \caption{1,3 Zoll OLED Display I2C SSH1106 Chip 128 x 64 Pixel}
        \label{fig:1_3_OLED_Display}
    \end{center}
\end{figure}

\section{Spezifikationen}
Entlang des Datenblatts, siehe \cite{azdelivery_oled_2021}, hat das Display die folgenden Spezifikationen:

\begin{description}
    \item[Modulabmessungen:] 36 mm x 34 mm x 3 mm
    \item[Punktmatrix:] 128 x 64 Pixel
    \item[Pixeldimension:] 0,21 mm × 0,21 mm
    \item[Pixelabstand:] 0,23 mm × 0,23 mm
    \item[Anzeigemodus:] Passive Matrix
    \item[Duty:] 1/64 Duty
    \item[Anzeigefarbe:] Weiß
    \item[Schnittstelle:] \ac{i2c}
    \item[Stromverbrauch:] Weniger als 11mA
    \item[Bildschirmgröße:] 1,3 Zoll
\end{description}

\subsubsection{Maximale Grenzwerte}
\begin{description}
    \item[Versorgungsspannung für Logik:] -0,3 V bis 5,5 V
    \item[Betriebstemperatur:] -40 °C bis +85 °C
    \item[Lagertemperatur:] -40 °C bis +85 °C
    \item[Lebensdauer bei 120 cd/m\textsuperscript{2}:] 10.000 Stunden
    \item[Lebensdauer bei 80 cd/m\textsuperscript{2}:] 30.000 Stunden
    \item[Lebensdauer bei 60 cd/m\textsuperscript{2}:] 50.000 Stunden
\end{description}

\subsubsection{Elektrische Eigenschaften}
\begin{description}
    \item[Versorgungsspannung für Logik:] 3,0 V bis 5,0 V (typisch 3,3 V)
    \item[Versorgungsspannung für Logik IO:] 3,0 V bis 3,3 V
    \item[Hoher Pegel Eingang:] 0,8 x VDDIO bis VDDIO
    \item[Niedriger Pegel Eingang:] 0 V bis 0,2 x VDDIO
    \item[Hoher Pegel Ausgang:] 0,8 x VDDIO bis VDDIO
    \item[Niedriger Pegel Ausgang:] 0 V bis 0,2 x VDDIO
    \item[Betriebsstrom für VDD:] 40 mA bis 45 mA
    \item[Schlafmodus Strom für VDD:] < 1 mA
\end{description}

\subsubsection{Anschlussbelegung}

Das Display verfügt über einen eingebauten 3,3V-Spannungsregler und kann entweder mit 3,3V oder 5V Logik und Stromversorgung betrieben werden, ohne das Display zu beschädigen.

\subsubsection{Anschlussdiagramme}

Das 1,3 Zoll \ac{oled} Display kann wie in Abbildung \ref{fig:SchaltplanOLED} an einen Mikrocontroller angeschlossen werden:

\begin{center}
    \includegraphics[width=12cm]{../../Images/OLED/OLEDCircuit04}
    \captionof{figure}{Schaltplan Displayverbindung mit Arduino}
    \label{fig:SchaltplanOLED}
\end{center}

\begin{description}
    \item[0/SCL:] A5
    \item[1/SDA:] A4
    \item[2/VDD:] 3,3V 
    \item[3/GND:] GND
\end{description}

\subsubsection{Vorsichtsmaßnahmen}
\begin{description}
    \item Die maximalen Grenzwerte für Logik- und LC-Treiber sind stets zu beachten.
    \item Eine umgekehrte Polarität von VDD und GND sollte vermieden werden.
    \item OLED-Module sollten nicht unter direkter Sonneneinstrahlung oder unter fluoreszierendem Licht gelagert werden.
    \item Mechanische Einwirkungen wie Fallenlassen, Biegen oder Verdrehen des Moduls sind zu vermeiden.
    \item Für die Lagerung sollten antistatische Behälter sowie eine saubere Umgebung verwendet werden.
    \item Zur Minimierung des „Screen Burn“-Effekts ist auf die Darstellung dauerhaft identischer Bilder zu verzichten.
\end{description}


\section{U8g2lib.h}
Im Folgenden wird auf die U8g2lib.h eingegangen. Es wird beschrieben, wofür diese Bibliothek zu gebrauchen ist, wie sie installiert werden kann, welche Funktionen sie bietet und ein Beispiel zur Anwendung gezeigt.

\subsection{Beschreibung}
Die U8g2lib.h-Bibliothek ist eine Bibliothek zur Ansteuerung von Monochrom-Grafikdisplays mit verschiedenen Schnittstellen wie I2C und SPI. Diese Bibliothek unterstützt eine Vielzahl von Display-Typen und ermöglicht die Erstellung von komplexen Grafiken und Texten auf kleinen \ac{oled}- und LCD-Displays.
Die Bibliothek U8g2lib.h stellt Funktionen zur Initialisierung und Ansteuerung von Grafikdisplays bereit. Mit \PYTHON{U8G2\_SSD1306\_128X64\_NONAME\_F\_HW\_I2C u8g2(U8G2\_R0, U8X8\_PIN\_NONE);} kann ein Display-Objekt erstellt werden, wobei der Display-Typ und die Schnittstelle spezifiziert werden. Die Initialisierung des Displays erfolgt durch den Aufruf von \PYTHON{u8g2.begin();}, wodurch das Display für die weitere Verwendung vorbereitet wird.
Um Grafiken und Texte auf dem Display anzuzeigen, bietet die Bibliothek eine Reihe von Funktionen zur Zeichnung von primitiven Formen wie Linien, Rechtecken, Kreisen und Polygonen. Zum Beispiel zeichnet \PYTHON{u8g2.drawLine(x0, y0, x1, y1);} eine Linie zwischen den Punkten (x0, y0) und (x1, y1), während \PYTHON{u8g2.drawBox(x, y, w, h);} ein gefülltes Rechteck mit der oberen linken Ecke bei (x, y) und den Abmessungen w und h erstellt.
Texte können mit \PYTHON{u8g2.drawStr(x, y, "Text");} auf das Display geschrieben werden, wobei (x, y) die Startkoordinaten des Textes sind. Die Schriftart lässt sich durch \PYTHON{u8g2.setFont(u8g2\_font\_ncenB08\_tr);} festlegen, was die ausgewählte Schriftart für alle nachfolgenden Textausgaben einstellt.

\subsection{Installation}
Die U8g2lib.h Bibliothek kann über den Library Manager der Arduino IDE installiert werden. Dafür sind folgende Schritte notwendig:

\begin{itemize}
    \item Arduino IDE öffnen
    \item In der Menüleiste:  \menu{Sketch>Bibliothek einbinden>Bibliotheken verwalten\ldots}
    \item Im Suchfeld \menu{u8g2} eingeben
    \item In Trefferliste die Bibliothek \menu{U8g2 by olikraus} auswählen
    \item Klicken auf \menu{Installieren}, um die Bibliothek zu der Entwicklungsumgebung hinzuzufügen.
\end{itemize}


\subsection{Funktionen}

\begin{itemize}
    \item \PYTHON{u8g2.clearBuffer()} - RAM-Zwischenspeicher leeren, während Display gleich bleibt
    \item \PYTHON{u8g2.drawStr()} - Textausgabe
    \item \PYTHON{u8g2.drawBox()} -  Gefülltes Rechteck
    \item \PYTHON{u8g2.drawCircle()} - Kreis zeichnen
    \item \PYTHON{u8g2.setFont()} - Schriftart setzen
    \item \PYTHON{u8g2.sendBuffer()} - Bufferinhalt ans Display senden
\end{itemize}



\subsection{Beispiel}
Der Beispielcode \ref{lst:u8g2} zeigt die Initialisierung eines SH1106 Treiber-basierten OLED-Displays über I2C. 

{
    \captionof{code}{Beispiel U8g2lib.h}\label{lst:u8g2}
    \ArduinoExternal{}{../../Code/Arduino/OLED/U8G2Example/U8G2Example.ino}
}


\section{Weiterführende Literatur}

Peter Adam Hoeher, Visible Light Communications \cite{hoeher_visible_2019}\\


Dieses Buch bietet einen umfassenden Überblick über die Theorie und Praxis der sichtbaren Lichtkommunikation (VLC). Es beschreibt die Nutzung von LED-Technologien zur gleichzeitigen Beleuchtung und Datenübertragung, wobei auch organische Leuchtdioden (OLEDs) als alternative Lichtquelle berücksichtigt werden. OLEDs werden als potenzielle Komponenten für zukünftige VLC-Systeme genannt, insbesondere aufgrund ihrer flexiblen Formfaktoren und flächigen Lichtabstrahlung, auch wenn ihre derzeitige Effizienz und Modulationsgeschwindigkeit noch hinter der von Standard-LEDs zurückbleibt.\\


Leonhard Stiny: Aktive elektronische Bauelemente \cite{Stiny:2019}\\



Dieses Buch stellt umfassendes Wissen zur Elektronikpraxis im Bereich der Halbleitertechnik bereit. Dabei werden Bauteile beschrieben, erläutert und erklärt, wie diese
in Schaltungen integriert werden können, um fehlerfreie und sichere Ergebnisse zu
erzielen. Auf zentrale Bauelemente wie unter anderem Dioden, bzw. LEDs (auch OLED) wird besonders Bezug genommen.